% =====================================================================
%  COMP3320 2026 -- Lab 1 introduction slides
%  Source of truth: script.md  (this file is the LaTeX translation)
%  Theme: cookie (LuaLaTeX-first).  Build with:  latexmk slides.tex
% =====================================================================
\documentclass[aspectratio=169]{beamer}

\usepackage{amsmath,amssymb}
\usepackage{booktabs}
\usepackage{tabularx}

% Logos etc. live in the shared template assets; add local assets/ too.
\graphicspath{{assets/}{../../template/assets/}}

\definecolor{fortranPurple}{HTML}{503874}     % poster deep purple
\definecolor{fortranPurpleMid}{HTML}{7E5AA6}
\usetheme[
  accent=fortranPurple,
  progressbar=foot,
  sectionpage=progressbar,
  subsectionpage=none,
  numbering=fraction,
  numberrail=section,
  block=fill,
  notes=hide,
  toc=aftertitle,
  closing=contact
]{cookie}

% small muted right-aligned annotation, as in the water deck
\newcommand{\cplx}[1]{{\footnotesize\color{cookieMuted}#1}}

\title[COMP3320 Lab 1]{How do I time my programs?}
\subtitle{Timing, scaling \& profiling in scientific HPC}
\author[COMP3320 2026 \textperiodcentered\ Lab 1]{Jorge Luis Gálvez Vallejo}
\cookieemail{jorge.galvezvallejo@anu.edu.au}
\institute{School of Computing}
\cookieuni{Australian National University}
\cookielocation{Canberra, Australia}
\date{Week 2, 2026}
\cookietagline{COMP3320 \textperiodcentered\ Lab 1 \textperiodcentered\ Week 2}

\cookieclosingtitle{}
\cookieclosingtext{Fork, clone, \texttt{make}, and run.}

\begin{document}

\maketitle

% =====================================================================
\section{Intro to scientific HPC}

\begin{frame}{Why should you care about HPC?}{It is the engine room of modern science and industry}
  \begin{columns}[T,onlytextwidth]
    \column{.5\textwidth}
      \textbf{Computational Fluid Dynamics}
      \begin{itemize}\setlength\itemsep{1pt}
        \item Aerospace --- Boeing, Lockheed, Airbus
        \item Cars --- F1, NASCAR
        \item Materials --- wind turbines, buildings
        \item Climate --- atmospheric \& oceanic physics
      \end{itemize}
      \smallskip
      \textbf{Computational Physics}
      \begin{itemize}\setlength\itemsep{1pt}
        \item Astro --- NASA, ESA, JPL
        \item High-energy --- nuclear, CERN, plasma
      \end{itemize}
      \smallskip
      \textbf{Computational Chemistry}
      \begin{itemize}\setlength\itemsep{1pt}
        \item Drug design, molecular dynamics --- big pharma
        \item Catalysts, materials --- energy, plastics
      \end{itemize}

    \column{.5\textwidth}
      \textbf{Software design \& optimization}
      \begin{itemize}\setlength\itemsep{1pt}
        \item HP-libraries --- NVIDIA, Intel, AMD, Cray
        \item Compilers --- Arm, NVIDIA, Intel, AMD, Cray
        \item Applied research, ML --- \emph{everywhere now}
      \end{itemize}
      \smallskip
      \textbf{Research \& collaboration}
      \begin{itemize}\setlength\itemsep{1pt}
        \item National labs all over the world
        \item HPC centres around the world
      \end{itemize}
      \vfill

  \end{columns}
\end{frame}

% =====================================================================
\section{How do I time my programs?}

\begin{frame}{Timers: two properties that matter}
  When you call a timing routine, ask two questions:
  \begin{itemize}\setlength\itemsep{6pt}
    \item \textbf{Resolution} --- how \emph{finely} can time be measured?
      \\ \cplx{milliseconds? microseconds? nanoseconds?}
    \item \textbf{Overhead} --- how long does it take to \emph{call} the timing
      function itself?
  \end{itemize}
  \bigskip
  \begin{block}{How you find out}
    Make repeated calls to the timer and look at the numbers it gives back.
  \end{block}
\end{frame}

\begin{frame}[fragile]{A timing loop --- what is it really measuring?}
  \begin{cookiecode}[title={the classic pattern}]{C}
double measure_time();
double elapsed_time(start, end);

for (int i = 0; i < ni; i++) {
  for (int j = 0; j < nj; j++) {
    for (int k = 0; k < nk; k++) {
      double t0 = measure_time();
      {
        do_ops(i, j, k);
      }
      double t1 = measure_time();
      double dt = elapsed_time(t0, t1);
      printf(" took = %f \n", dt);
    }
  }
}
  \end{cookiecode}
\end{frame}

\begin{frame}{Timing}
  \begin{itemize}\setlength\itemsep{6pt}
    \item The previous program is a 3-fold loop: \texttt{measure\_time} is called
      $n_i \times n_j \times n_k$ times.
    \item If \texttt{do\_ops(i,j,k)} takes ${\sim}1\,\text{ms}$ \textbf{and the
      timing routine has a similar overhead}\ldots
  \end{itemize}
  \bigskip
  \begin{cookiequote}
    What are you then measuring?
  \end{cookiequote}
  \bigskip
  \cplx{How accurate is the \texttt{measure\_time} routine in the first place?}
\end{frame}

\begin{frame}{Resolution}
  \begin{itemize}\setlength\itemsep{8pt}
    \item If my clock can only measure \textbf{milliseconds} accurately, what
      happens when my routines take \textbf{microseconds}?
      \\ \cplx{the clock reports 0, or noise}
    \item Timing routines vary across \textbf{languages, libraries, and operating
      systems}.
  \end{itemize}
  \bigskip
  \begin{block}{Rule of thumb}
    Match the clock to the thing you are timing.
  \end{block}
\end{frame}

% =====================================================================
\section{Benchmarking \& profiling}

\begin{frame}{Benchmarking and appropriate timing routines}
  \begin{columns}[T,onlytextwidth]
    \column{.5\textwidth}
      \textbf{Benchmarking} --- time it from the outside
      \begin{itemize}\setlength\itemsep{3pt}
        \item If you are timing wrong, you are \emph{not} benchmarking properly.
        \item Proper benchmarking relies on a \textbf{good clock} used correctly.
      \end{itemize}
    \column{.5\textwidth}
      \textbf{Profiling} --- look inside the running program
      \begin{itemize}\setlength\itemsep{3pt}
        \item Instrument the code at the \textbf{binary level} to extract information.
        \item Use instruction- and runtime-analysis for precise measurements.
        \item Tells you \emph{where} the time actually goes --- often somewhere surprising.
      \end{itemize}
  \end{columns}
\end{frame}

% =====================================================================
\section{Lab 1: iterative non-linear solvers}

\begin{frame}{The program you will study: the Bratu problem}
  A \emph{nonlinear} PDE, the \textbf{Liouville--Bratu--Gelfand equation}:
  \[
    -\nabla^2 u = \lambda\, e^{u}
  \]
  \begin{itemize}\setlength\itemsep{4pt}
    \item Solved on the unit square, with $u = 0$ on the boundary\ldots
    \item \ldots discretised onto an $n \times n$ grid.
  \end{itemize}
\end{frame}

\begin{frame}{Where does this equation come from?}{Optional colour --- it makes $\lambda$ less mysterious}
  A classic model of \textbf{thermal self-ignition} --- a balance of two terms:
  \begin{itemize}\setlength\itemsep{4pt}
    \item $-\nabla^2 u$: how fast heat \textbf{conducts away}.
    \item $e^{u}$: how fast heat is \textbf{generated}
      --- hotter $\Rightarrow$ reacts faster $\Rightarrow$ a feedback loop.
    \item $\lambda$ sets generation vs.\ conduction.
  \end{itemize}
  \bigskip
  \begin{alertblock}{The interesting physics}
    Above a critical $\lambda^{*} \approx 6.808$ \textbf{no steady state exists} ---
    the material cannot shed heat fast enough and it ignites.
    We fix $\lambda = 3.0$, safely below critical, so a solution exists.
  \end{alertblock}
\end{frame}

\begin{frame}{How the code solves it}{You do not need the details --- just appreciate the shape}
  \begin{itemize}\setlength\itemsep{3pt}
    \item Nonlinear $\Rightarrow$ we cannot solve it in one step like $Ax = b$.
    \item Repeatedly apply a fixed-point map $u \leftarrow G(u)$ (one grid sweep),
      accelerated by \textbf{DIIS} (Anderson acceleration).
  \end{itemize}
  \medskip
  \begin{center}
    \begin{tabular}{@{}ll@{}}
      \toprule
      \textbf{phase} & \textbf{what it does} \\
      \midrule
      \texttt{map (G)}     & one sweep of the grid \\
      \texttt{error vec}   & forms $e = G(x) - x$ \\
      \texttt{B matrix}    & dot products of stored error vectors \\
      \texttt{small solve} & tiny dense solve for the DIIS coefficients \\
      \texttt{extrapolate} & combines history into the next iterate \\
      \bottomrule
    \end{tabular}
  \end{center}
  \medskip
  \cplx{The code reports the time spent in each phase --- that is your raw data.}
\end{frame}

\begin{frame}{What you'll do today}
  \begin{enumerate}\setlength\itemsep{6pt}
    \item \textbf{Timer resolution \& overhead} --- \texttt{gettimeofday} vs
      \texttt{clock\_gettime}.
    \item \textbf{CPU timing \& scaling} --- how does one iteration scale,
      $\mathcal{O}(n^{k})$? Does a better \emph{algorithm} (DIIS) beat brute force?
    \item \textbf{Profiling with Intel tools} --- VTune (hotspots) \& Advisor (loops).
    \item \textbf{Profiling with Linaro Forge} --- a second profiler; see what the
      \emph{compiler} did.
  \end{enumerate}
\end{frame}

\begin{frame}[standout,noframenumbering]
  Run the code. 
\end{frame}

\begin{frame}{The golden rule}
  \begin{itemize}\setlength\itemsep{6pt}
    \item The compiler and the hardware do things you \textbf{cannot see} by
      reading \texttt{diis.c}.
    \item Almost every answer comes from \textbf{measurement on Gadi}, not
      inspection.
  \end{itemize}
\end{frame}

\begin{frame}[fragile]{Getting started}
  \begin{enumerate}\setlength\itemsep{4pt}
    \item \textbf{Fork} the repo to your own account (as in the pre-lab).
    \item \textbf{Clone} it on Gadi:
  \end{enumerate}
  \begin{cookiecode}[title={clone}]{bash}
git clone git@gitlab.comp.anu.edu.au:<YOUR_UID>/comp3320-2026-lab-1.git
  \end{cookiecode}
  \begin{enumerate}\setlength\itemsep{4pt}\setcounter{enumi}{2}
    \item \textbf{Build \& run}:
  \end{enumerate}
  \begin{cookiecode}[title={build}]{bash}
make walltime      # Part 1
make exe_diis      # Parts 2-4
./exe_diis 64
  \end{cookiecode}
\end{frame}

% closing slide is inserted automatically by `closing=contact`
\end{document}

